
\section{Discussion}
\label{sec:discussion}

\paragraph{Our Contributions.}
Our results highlight a major weakness in prior ASCA approaches -- their vulnerability to noise. While previous methods achieved high accuracy in noise-free conditions, their performance deteriorates in real-world settings. We address this gap by leveraging transformer-based architectures (VTs and LLMs) for robust classification and noise correction. Our findings underscore the need for ASCA research to prioritize noise resilience and adopt transformer-driven methodologies.
% Our results underscore a critical yet often overlooked weakness of prior ASCA approaches: their substantial vulnerability in realistic, noisy conditions. Previous methodologies—ranging from statistical and traditional machine learning techniques to more sophisticated deep-learning models (CNNs and CoAtNet)—typically reported high accuracies only under ideal (noise-free) scenarios. However, as we demonstrate, their effectiveness quickly deteriorates when confronted with realistic ambient noise. This limitation highlights an essential reality-check for future ASCA research; evaluating methodologies solely on clean datasets may result in an overly optimistic understanding of their real-world efficacy. In this study, we proposed an effective and relatively unexplored solution -- leveraging transformer-based architectures (particularly VTs and LLMs) -- to remedy this critical issue. Our empirical validation demonstrates transformers' strong potential for robust classification and their unique ability to mitigate noise-induced errors through global correlation capabilities and contextual understanding. Future ASCA research should thus treat noise robustness seriously, and methodologies incorporating transformers and LLM-driven contextual correction should be considered imperative.

\paragraph{Implications of Findings.}
Our findings challenge the assumption that passphrases provide both security and usability against LLM-assisted ASCA attacks. While prior research shows passphrases are easier to remember and more secure than simple passwords \cite{keith2009behavioral}, LLMs can reconstruct them using linguistic patterns. However, they struggle with entirely random inputs, which are impractical for users. As transformer-based ASCAs grow more effective, both passwords and passphrases become unreliable. This reinforces the need to shift towards alternative authentication methods like hardware-based security, Multi-Factor Authentication (MFA), or behavioral biometrics, which are more resistant to side-channel inference.
% Our findings suggest that the long-standing assumption that passphrases offer better security while maintaining usability may no longer hold in the presence of LLM-assisted ASCA attacks. Prior research has shown that passphrases are easier to remember while providing stronger security than simple passwords~\cite{keith2009behavioral}, but our results indicate that LLMs can effectively reconstruct structured passphrases by leveraging their inherent linguistic patterns. Conversely, LLM-based correction fails when no discernible pattern exists, which would require users to adopt completely unpredictable, random inputs—an approach that is highly impractical from a usability perspective. This highlights a fundamental issue with password-based authentication: as ASCA attacks become more powerful with transformer-based models, neither passwords nor passphrases remain reliable defenses. Given these findings, our work reinforces the argument that password-based authentication is becoming obsolete, and future security efforts should focus on alternative authentication mechanisms such as hardware-based authentication, Multi-Factor Authentication (MFA), or behavioral biometrics, which are less susceptible to side-channel inference.

\paragraph{This Work's Limitations.}
Our study has limitations, particularly in dataset size and scope. It includes only 25 samples per keystroke, limited to alphanumeric inputs, excluding essential keys like space, backspace, and enter. Recordings were restricted to a MacBook Pro keyboard, limiting generalizability across devices and typing styles. Additionally, we used synthetic Gaussian noise, whereas real-world environments involve more complex noise sources. While we mitigated these factors by varying noise levels and using standard keyboards, our key contribution lies in demonstrating the feasibility of transformer-based contextual error correction, paving the way for future noise-resilient ASCA research.
% Despite these important results, our work, like any scientific investigation, has limitations that must be acknowledged. The dataset used in this study is relatively limited, both in terms of size and scope. Specifically, it consists of only 25 samples per keystroke, limited to standard alphanumeric inputs (letters and digits). Essential keyboard keys widely used in realistic typing scenarios—such as spaces, backspace, enter, arrows, function, and special symbol keys—were not included in our dataset. Furthermore, recordings were restricted to a single keyboard type (MacBook Pro keyboard); consequently, results might not generalize directly across diverse device models or user typing behaviors. Additionally, the acoustic noise model employed synthetic Gaussian noise, whereas real-world acoustic environments involve complex sources like background conversations, music, and intermittent ambient disturbances. Although these factors represent limitations, we partially mitigated their effects by systematically varying noise intensities and selecting standard, widely-used keyboard environments. More critically, our study's key contribution—demonstrating the feasibility and effectiveness of transformer-based contextual error correction—offers value beyond isolated numerical results by establishing a broader, conceptual paradigm for future noise-resilient research in ASCAs.

\paragraph{This Fields's Limitations.}
ASCA research lacks large, diverse, and public datasets, limiting progress and reproducibility. Unlike speech recognition or NLP, it relies on small, non-standardized datasets, making benchmarking difficult. The absence of varied keystroke-acoustic data hinders robust model development. To advance the field, community-driven dataset collection is crucial. Without standardized, large-scale datasets, progress will remain fragmented, and security risks harder to assess. We urge collaborative efforts to establish open ASCA benchmarks, similar to those in speech and image recognition, to drive systematic advancements.
% A significant limitation in ASCA research is the lack of large, diverse, and publicly available datasets, which severely constrains progress and reproducibility in the field. Unlike speech recognition or NLP, where large-scale, well-curated datasets exist, ASCA studies often rely on small, non-standardized, and proprietary datasets, making it difficult to benchmark new methods against prior work. The absence of widely available keystroke-acoustic datasets covering varied keyboard types, ambient noise conditions, and diverse user typing behaviors hinders the development of more robust and generalizable models. Moving forward, the field must prioritize community-driven dataset collection efforts, enabling researchers to build, test, and refine ASCA techniques under standardized conditions. Without a publicly available, large-scale ASCA dataset, progress will remain fragmented, and the broader security implications of these attacks will be difficult to assess comprehensively. We call for collaborative efforts to establish open ASCA benchmarks, much like those in speech and image recognition, to drive systematic advancements in the field.

\paragraph{Our Future Work.}
Given these limitations, future research has significant potential. Expanded datasets should include additional keyboard keys (function, space, backspace, arrows, special characters), diverse device models (desktop and soft keyboards on smartphones/tablets), and varied typing habits to improve generalization. Critically, experiments should incorporate authentic ambient noise (e.g., coffee shop chatter, street noise, music, office sounds) instead of synthetic models alone. While our study introduced realism by adjusting noise levels, future work should explicitly simulate diverse acoustic scenarios and thoroughly evaluate performance. Additionally, exploring real-time error correction with lightweight, fine-tuned LLMs can enhance robust side-channel methods on constrained edge devices, requiring efficient inference, low latency, and improved usability. We encourage prioritizing these directions to advance practical ASCA solutions.
% Given these limitations, significant opportunities exist for future research. Expanded and diversified databases should incorporate additional keyboard keys (including function, space, backspace, arrows, special characters), varied device models (desktop keyboards, soft keyboards on smartphones/tablets), and different user typing habits to ensure robust generalizations. Critically, future experimental methodologies must integrate authentic ambient auditory noise, such as coffee shop background chatter, street noises, music, and office sounds, rather than relying on synthetic noise models alone. While our study took an essential initial step towards realism by systematically adjusting noise levels, future work should explicitly simulate varied realistic acoustic scenarios and comprehensively evaluate model performance under each condition. Furthermore, real-time error correction with lightweight, fine-tuned LLMs remains an exciting research direction, with potential applications for deploying robust side-channel recovery methods on computationally constrained edge devices. Such experiments would necessitate optimizing inference efficiency, minimizing latency, and improving usability. We encourage future research to prioritize these directions to advance practical solutions to real-world ASCA challenges.

\paragraph{Field's Future Work.}
Future ASCA research will extend transformer-based methods beyond acoustics into visual side-channels (screen reflections, typing videos) and electromagnetic emanations. As LLMs can reconstruct passphrases despite noise, password-based authentication may become less viable, prompting shifts toward multi-factor authentication, behavioral biometrics, or continuous authentication. Addressing the lack of large, standardized public datasets covering diverse keyboards, typing styles, and real-world noise will be critical for reproducibility. Finally, efficient transformer optimizations (LoRA, QLoRA, distillation) could enable real-time, on-device ASCA attacks, increasing security risks and driving demand for new cybersecurity countermeasures.

% The future of ASCA research will likely see transformer-based methodologies applied beyond acoustic attacks, extending to visual side-channels, such as screen reflections, typing video analysis, and even electromagnetic emanations. With LLMs demonstrating the ability to reconstruct passphrases despite noise, this raises concerns about the long-term viability of password-based authentication, potentially accelerating a shift towards multi-factor authentication, behavioral biometrics, or continuous authentication mechanisms. Another pressing challenge is the lack of large, standardized public datasets, limiting comparability and reproducibility across studies; future work must prioritize community-driven dataset collection covering diverse keyboards, typing styles, and real-world noise conditions. Additionally, as transformer architectures become more efficient through optimizations like LoRA, QLoRA, and model distillation, we expect to see real-time, on-device ASCA attacks become feasible, further increasing security risks and necessitating new countermeasures in cybersecurity defenses.